\documentclass[12pt]{article}
\usepackage{amsmath, amssymb}
\usepackage[margin=1in]{geometry}
\setlength{\parskip}{6pt}
\title{QUBO Formulation: Minimum Vertex Cover}
\date{}
\begin{document}
\maketitle

\subsection*{Minimum Vertex Cover}

\paragraph{Problem Statement.}
Given an undirected graph $G = (V, E)$ with vertex set $V$ and edge set $E$, the Minimum Vertex Cover problem seeks a subset $S \subseteq V$ of minimum cardinality such that every edge in $E$ is incident to at least one vertex in $S$. Formally, we wish to minimize $|S|$ subject to the constraint that for every edge $(u, v) \in E$, the intersection $\{u, v\} \cap S$ is non-empty.

\paragraph{Decision Variables.}
We introduce a binary variable $x_i \in \{0, 1\}$ for each vertex $i \in V$. The variable $x_i$ indicates membership in the cover:
\[
x_i = \begin{cases} 
1 & \text{if } i \in S, \\
0 & \text{if } i \notin S.
\end{cases}
\]
The size of the cover is given by $\sum_{i \in V} x_i$.

\paragraph{QUBO Derivation.}
We construct a Quadratic Unconstrained Binary Optimization (QUBO) Hamiltonian $H(x)$ whose global minimum over $\{0,1\}^{|V|}$ corresponds to the optimal vertex cover.

\textbf{Step 1: Objective Function.}
The objective is to minimize the number of selected vertices:
\[
\min \sum_{i \in V} x_i.
\]
In the QUBO formulation, this contributes directly to the Hamiltonian as a linear term with positive coefficients:
\begin{align}
H_{\text{obj}}(x) = \sum_{i \in V} x_i.
\end{align}

\textbf{Step 2: Constraints.}
For each edge $(u, v) \in E$, the vertex cover constraint requires that at least one endpoint is selected:
\[
x_u + x_v \ge 1, \quad \forall (u, v) \in E.
\]
This constraint is violated if and only if $x_u = 0$ and $x_v = 0$.

\textbf{Step 3: Penalty Terms.}
To enforce the constraints, we add a penalty term for each edge that is zero when the constraint is satisfied and positive otherwise. A suitable penalty for the constraint $x_u + x_v \ge 1$ is:
\[
P_{uv}(x) = \lambda (1 - x_u)(1 - x_v),
\]
where $\lambda > 0$ is a penalty weight. Expanding this expression yields:
\begin{align}
P_{uv}(x) &= \lambda (1 - x_u - x_v + x_u x_v) \nonumber \\
&= \lambda - \lambda x_u - \lambda x_v + \lambda x_u x_v.
\end{align}
Note that since $x_i \in \{0,1\}$, we have $x_i^2 = x_i$, but no quadratic terms arise in the expansion of $(1-x_u)(1-x_v)$. The penalty is $\lambda$ when $x_u=x_v=0$ and $0$ otherwise, correctly penalizing only infeasible assignments.

\textbf{Step 4: Combined Hamiltonian.}
The total Hamiltonian is the sum of the objective and all edge penalties:
\begin{align}
H(x) &= \sum_{i \in V} x_i + \sum_{(u, v) \in E} \left( \lambda - \lambda x_u - \lambda x_v + \lambda x_u x_v \right).
\end{align}
We collect terms by grouping the constant offset, linear terms, and quadratic terms. The constant offset is:
\begin{align}
c = \sum_{(u, v) \in E} \lambda = \lambda |E|.
\end{align}
The coefficient of a linear variable $x_i$ is $1$ from the objective plus contributions from each incident edge. Let $\text{deg}(i)$ denote the degree of vertex $i$ in $G$. Each incident edge $(i, j)$ contributes $-\lambda x_i$. Thus:
\begin{align}
\text{Coeff}(x_i) = 1 - \lambda \cdot \text{deg}(i).
\end{align}
The coefficient of a quadratic term $x_u x_v$ is $\lambda$ if $(u, v) \in E$, and $0$ otherwise.

\textbf{Step 5: QUBO Matrix Entries.}
The Hamiltonian can be written in standard QUBO form $H(x) = c + \sum_{i \in V} Q_{i,i} x_i + \sum_{(u, v) \in E} Q_{u,v} x_u x_v$. The entries of the upper-triangular Q matrix and the offset are given by:
\begin{align}
Q_{i,i} &= 1 - \lambda \cdot \text{deg}(i), \quad \forall i \in V, \\
Q_{u,v} &= \lambda, \quad \forall (u, v) \in E, \\
Q_{u,v} &= 0, \quad \forall (u, v) \notin E, \\
c &= \lambda |E|.
\end{align}

\paragraph{Penalty Weight Justification.}
The penalty weight $\lambda$ must be chosen large enough to ensure that any infeasible solution has strictly higher energy than the best feasible solution. Consider a feasible solution $x^{\text{feas}}$ and an infeasible solution $x^{\text{infeas}}$. The maximum reduction in the objective function achievable by moving from a feasible to an infeasible state is bounded by the maximum possible objective value, which is at most $|V|$. However, a tighter bound arises from local violations: violating a single constraint allows setting both endpoints to $0$, reducing the objective by at most $2$ (if both were $1$) or $1$ (if one was $1$). More generally, for an objective with unit weights, the maximum gain from violating one constraint is $1$ per variable involved. A sufficient condition is $\lambda > \max_{i} |w_i|$, where $w_i$ are the objective coefficients. Here $w_i = 1$, so $\lambda > 1$ suffices. In practice, $\lambda = 2$ is chosen to provide a margin, ensuring that the energy cost of any violation dominates any potential objective improvement.

\paragraph{Correctness Argument.}
Minimizing $H(x)$ over $\{0,1\}^{|V|}$ recovers the optimal vertex cover. For any feasible solution, all penalty terms vanish, so $H(x) = \sum_{i \in V} x_i$, which equals the cover size. For any infeasible solution, at least one penalty term contributes $\lambda$ to the energy. Since $\lambda > 1$, the energy of any infeasible solution is strictly greater than that of any feasible solution with cover size less than $\lambda$. Specifically, if $x$ is infeasible, there exists a feasible solution $x'$ obtained by flipping a $0$ to $1$ in a violated edge; this reduces the penalty by at least $\lambda$ and increases the objective by at most $1$. Since $\lambda > 1$, the energy decreases, implying the global minimum cannot be infeasible. Among feasible solutions, $H(x)$ equals the objective function, so the minimum corresponds to the minimum vertex cover.

\paragraph{Illustrative Example.}
Consider a path graph with $N=3$ vertices $V=\{0, 1, 2\}$ and edges $E=\{(0, 1), (1, 2)\}$. Here $|E|=2$ and degrees are $\text{deg}(0)=1$, $\text{deg}(1)=2$, $\text{deg}(2)=1$. We set the penalty weight $\lambda = 2$.

Using the general formulas:
\begin{align}
c &= 2 \times 2 = 4, \\
Q_{0,0} &= 1 - 2(1) = -1, \\
Q_{1,1} &= 1 - 2(2) = -3, \\
Q_{2,2} &= 1 - 2(1) = -1, \\
Q_{0,1} &= 2, \quad Q_{1,2} = 2.
\end{align}
The resulting Hamiltonian is:
\begin{align}
H(x) = 4 - x_0 - 3x_1 - x_2 + 2x_0 x_1 + 2x_1 x_2.
\end{align}
Evaluating $H(x)$ for candidate solutions:
\begin{itemize}
    \item $x = (0, 1, 0)$ (cover $\{1\}$, size 1): $H = 4 - 0 - 3 - 0 + 0 + 0 = 1$.
    \item $x = (1, 0, 1)$ (cover $\{0, 2\}$, size 2): $H = 4 - 1 - 0 - 1 + 0 + 0 = 2$.
    \item $x = (0, 0, 0)$ (infeasible): $H = 4$.
    \item $x = (1, 1, 1)$ (cover, size 3): $H = 4 - 1 - 3 - 1 + 2 + 2 = 3$.
\end{itemize}
The global minimum is $H=1$ at $x=(0, 1, 0)$, corresponding to the minimum vertex cover $\{1\}$.
\end{document}